\scanpage{76}
\mspage{25}

Revenons à la prop. 2. La stabilité de la notion de morphisme asphérique par composition \diagram{\begin{tikzcd}[ampersand replacement=\&,row sep=small] X\arrow[d,"f"']\&X'\arrow[l,"r"']\arrow[d,"f'"]\\Y\arrow[d,"g"']\&Y'\arrow[l,"q"']\arrow[d,"{g'}"]\\ Z\&Z'\arrow[l,"p"]\end{tikzcd}} résulte aussitôt de la forme a) de la condition. De même pour la condition que si $gf$ et $f$ sont $W$-asphériques, $g$ aussi -- il faut montrer que si $Z'\longrightarrow Z$ est iso loc., alors $g'\deleted{\unsure{f'}}$ est dans $W$, or l'hyp. implique que $g'f'$ et $f'$ le sont.

\deleted{\uline{NB} \unsure{Quand}} \uline{Corollaire 1} L'image inverse d'un $f:X\longrightarrow Y$ qui est $W$-asphérique \inserted{(resp. $W$-coasphérique)} par un iso local (resp. par un co-iso local) est $W$-asphérique ($W$-coasphérique)

\diagram{\begin{tikzcd}[ampersand replacement=\&,column sep=small] X\arrow[d,dash]\&X'\arrow[l,dash]\arrow[d,dash]\&X''\arrow[l]\arrow[d]\\ \unsure{Y}\&Y'\arrow[l,"{\mathrm{iso\ loc}}"]\&Y''\arrow[l,"{\mathrm{iso\ loc}}"]\end{tikzcd}} Résulte formellement de la définition par a), et du fait qu'un composé de deux iso. locaux (resp coisom locaux) est itou.

\uline{NB} Bien sûr, la notion de foncteur $W$-asphérique n'est pas stable par chgt. de base quelconque. \deleted{Mais} Par exemple, l'inclusion $X=\{y_0\}\mathrel{\relbar}Y$ est $W$-asph. \unsure{sss|ssi} les catégories discrètes $\mathrm{Hom}\,(y_0,y)$ sont $W$-asph. i.e. (sous réserve \unsure{que l'on ait|qu'elles ont} $W(5)$, axiome de connexité) sss les $\mathrm{Hom}(y_0,y)$ sont de cardinal 1, i.e. $y_0$ \uline{objet initial} de $Y$. Mais si $Y\ne\{y_0\}$, prenant $y$ dans $Y$, $y\ne y_0$, \deleted{la fibre de} et prenant $Y'=\{y\}\longrightarrow Y$, \unsure{base|bas.} $X'=X\times_Y Y'=\varnothing$, et $\varnothing\longrightarrow e$ n'est pas asphérique !

% archive sheet number: 76
